\section{Observability and evaluation}
\label{sec:eval}

\subsection{Tracing}
Every check writes a JSON-lines trace under \texttt{data/traces/}, one file per check. Each line carries a timestamp, the check identifier, the event (step, question, answer, learning, verdict or error), the action, the planner source (rules, LLM, or rules after a fallback), a one-line summary, the duration and, for model calls, the model name, token counts and whether a fallback occurred. The same events stream to the client, which draws them as a live timeline next to the batch, shown as a strip of character cells that turn from \emph{read} to \emph{uncertain} or \emph{confirmed}. Writing a trace can never fail a check: errors in the tracer are swallowed, and the file is bounded when it is read back. Spans can additionally be exported through OpenTelemetry to a local Arize Phoenix server using the OpenInference conventions \cite{otel2026genai,arize2026openinference}; the export is switched on by a setting, probes the collector with a short timeout, and a check runs the same without it. Each check becomes one trace: an agent span for the check, a chain span per planner decision with its tool or model span beneath, and a guardrail span for the verifier.

\subsection{Test sets}
\emph{Synthetic diagnostics.} Eleven authored scenarios, each with an expected verdict, cover a clean listed batch, a batch absent from the index, a one-character misread, a spurious entry, a missing carton, a batch collision, missing product text, an unreadable batch, a date disagreement, a neighbouring sibling batch and an unreadable pack. They run offline in the test suite and from the interface, and they check behaviour, not accuracy on real medicines.

\emph{Evaluation set.} \res{ResN} cartons are rendered from real alert rows in the local index, each marked as a demonstration prop, with a fixed scenario mix: \res{ResMixClean} clean, \res{ResMixQr} with a QR code, \res{ResMixGlare} with glare, \res{ResMixMisread} misread, \res{ResMixSpurious} spurious, \res{ResMixCollision} collision, \res{ResMixGenuine} unlisted and \res{ResMixSibling} sibling packs. Thresholds were tuned on a development set (seed 1) only, and results are reported on a separate evaluation set (seed \res{ResSeed}). Collision packs carry a listed batch number on a different product and maker, which models the reuse of short codes seen in the data. Misread and sibling packs print the misread string on the carton, so the OCR error is simulated and the oracle answers from the batch the pack was meant to carry. The evaluation seed was run twice, with three code fixes in between; Section~\ref{sec:runs} reports both runs.

\subsection{Systems compared}
The baseline is a lookup script: QR code or OCR, a regular expression for the batch, exact lookup, and any hit reported as an alarm. It is compared with the agent using the rule planner and the agent using the hybrid planner. Each arm runs twice: with an oracle user that answers every question truthfully from the ground truth, and with every question skipped, which shows what the questions buy.

\subsection{Metrics and statistics}
Hit recall is the share of packs whose truth is \FLAGGED{} or \SFAKE{} that receive the correct label. Type accuracy is the share of detected positives with the right one of the two labels. The false-alarm rate is the share of packs whose truth is \NLISTED{} that receive an alarm, reported overall and on collision packs alone. We also report the abstention rate, questions per pack, median and 95th-percentile latency, and safety violations (banned phrases or alarms without a citation, required to be zero). Agent comparisons without cost are misleading \cite{kapoor2024agents}, so latency and questions are reported for every arm. Rates carry 95\% Wilson intervals \cite{wilson1927probable,brown2001interval}, and paired comparisons on the same packs use the exact McNemar test \cite{mcnemar1947note,dietterich1998approximate}. With at most \res{ResMixClean} packs per scenario a difference of one or two packs is noise, so per-scenario results are given as counts.

\subsection{Results}
\label{sec:results}
\ifresults
The tables in this section give the headline metrics with the oracle user, the same metrics with every question skipped and the per-scenario counts. They are from the second run, after the fixes of Section~\ref{sec:runs}. Perception is shared by all arms and reads each image once; the batch read equalled the rendered text on \res{ResOCRBatchExact} images, at a median of \res{ResOCRMedian}\,s per image. All numbers are written by the evaluation command into \texttt{paper/generated/} and are not edited by hand.
\IfFileExists{generated/results_tables.tex}{% Generated by cli.py eval. Do not edit by hand. Needs booktabs and graphicx.
\begin{table}[t]\centering\footnotesize
\caption{Headline metrics on the eval set (seed 7, $n=60$), HITL: oracle user. Rates as value ($k/n$; Wilson 95\% interval).}
\label{tab:results-oracle}
\resizebox{\columnwidth}{!}{%
\begin{tabular}{llll}
\toprule
Metric & Baseline (lookup) & Agent (rules) & Agent (hybrid) \\
\midrule
Hit recall & 0.80 (32/40; 0.65-0.90) & 0.82 (33/40; 0.68-0.91) & 0.82 (33/40; 0.68-0.91) \\
Hit recall (clean + QR) & 1.00 (16/16; 0.81-1.00) & 0.75 (12/16; 0.51-0.90) & 0.75 (12/16; 0.51-0.90) \\
Type accuracy & 1.00 (32/32; 0.89-1.00) & 1.00 (33/33; 0.90-1.00) & 1.00 (33/33; 0.90-1.00) \\
False-alarm rate (all NOT\_LISTED) & 0.60 (12/20; 0.39-0.78) & 0.00 (0/20; 0.00-0.16) & 0.00 (0/20; 0.00-0.16) \\
False-alarm rate (collisions) & 1.00 (8/8; 0.68-1.00) & 0.00 (0/8; 0.00-0.32) & 0.00 (0/8; 0.00-0.32) \\
False-alarm rate (siblings) & 1.00 (4/4; 0.51-1.00) & 0.00 (0/4; 0.00-0.49) & 0.00 (0/4; 0.00-0.49) \\
Abstention rate & 0.00 (0/60; 0.00-0.06) & 0.00 (0/60; 0.00-0.06) & 0.00 (0/60; 0.00-0.06) \\
Overall accuracy & 0.67 (40/60; 0.54-0.77) & 0.88 (53/60; 0.78-0.94) & 0.88 (53/60; 0.78-0.94) \\
Questions per pack & 0.00 & 0.55 & 0.55 \\
Steps per pack & 1.0 & 7.8 & 7.8 \\
Latency median / mean / p95 (s) & 1.33 / 1.36 / 1.82 & 1.44 / 1.48 / 1.89 & 1.51 / 3.52 / 9.64 \\
Reasoning only, median / p95 (s) & 0.000 / 0.000 & 0.125 / 0.208 & 0.097 / 8.269 \\
Safety violations & 0 & 0 & 0 \\
LLM steps / fallback steps & 0 / 0 & 0 / 0 & 27 / 0 \\
\bottomrule
\end{tabular}}
\end{table}

\begin{table}[t]\centering\footnotesize
\caption{Headline metrics on the eval set (seed 7, $n=60$), HITL: every question skipped. Rates as value ($k/n$; Wilson 95\% interval).}
\label{tab:results-none}
\resizebox{\columnwidth}{!}{%
\begin{tabular}{llll}
\toprule
Metric & Baseline (lookup) & Agent (rules) & Agent (hybrid) \\
\midrule
Hit recall & 0.80 (32/40; 0.65-0.90) & 0.12 (5/40; 0.05-0.26) & 0.12 (5/40; 0.05-0.26) \\
Hit recall (clean + QR) & 1.00 (16/16; 0.81-1.00) & 0.31 (5/16; 0.14-0.56) & 0.31 (5/16; 0.14-0.56) \\
Type accuracy & 1.00 (32/32; 0.89-1.00) & 1.00 (5/5; 0.57-1.00) & 1.00 (5/5; 0.57-1.00) \\
False-alarm rate (all NOT\_LISTED) & 0.60 (12/20; 0.39-0.78) & 0.00 (0/20; 0.00-0.16) & 0.00 (0/20; 0.00-0.16) \\
False-alarm rate (collisions) & 1.00 (8/8; 0.68-1.00) & 0.00 (0/8; 0.00-0.32) & 0.00 (0/8; 0.00-0.32) \\
False-alarm rate (siblings) & 1.00 (4/4; 0.51-1.00) & 0.00 (0/4; 0.00-0.49) & 0.00 (0/4; 0.00-0.49) \\
Abstention rate & 0.00 (0/60; 0.00-0.06) & 0.52 (31/60; 0.39-0.64) & 0.52 (31/60; 0.39-0.64) \\
Overall accuracy & 0.67 (40/60; 0.54-0.77) & 0.37 (22/60; 0.26-0.49) & 0.37 (22/60; 0.26-0.49) \\
Questions per pack & 0.00 & 0.52 & 0.52 \\
Steps per pack & 1.0 & 7.6 & 7.6 \\
Latency median / mean / p95 (s) & 1.33 / 1.36 / 1.82 & 1.45 / 1.49 / 1.91 & 1.52 / 3.32 / 12.78 \\
Reasoning only, median / p95 (s) & 0.000 / 0.000 & 0.132 / 0.190 & 0.115 / 11.363 \\
Safety violations & 0 & 0 & 0 \\
LLM steps / fallback steps & 0 / 0 & 0 / 0 & 30 / 0 \\
\bottomrule
\end{tabular}}
\end{table}

\begin{table}[t]\centering\footnotesize
\caption{Correct verdicts per scenario ($k/n$), HITL: oracle.}
\label{tab:scenario-oracle}
\begin{tabular}{lccc}
\toprule
Scenario & Baseline (lookup) & Agent (rules) & Agent (hybrid) \\
\midrule
clean & 8/8 & 7/8 & 7/8 \\
glare & 8/8 & 7/8 & 7/8 \\
misread & 0/8 & 6/8 & 6/8 \\
qr & 8/8 & 5/8 & 5/8 \\
spurious & 8/8 & 8/8 & 8/8 \\
collision & 0/8 & 8/8 & 8/8 \\
genuine & 8/8 & 8/8 & 8/8 \\
sibling & 0/4 & 4/4 & 4/4 \\
\bottomrule
\end{tabular}
\end{table}

\begin{table}[t]\centering\footnotesize
\caption{Correct verdicts per scenario ($k/n$), HITL: questions skipped.}
\label{tab:scenario-none}
\begin{tabular}{lccc}
\toprule
Scenario & Baseline (lookup) & Agent (rules) & Agent (hybrid) \\
\midrule
clean & 8/8 & 0/8 & 0/8 \\
glare & 8/8 & 0/8 & 0/8 \\
misread & 0/8 & 0/8 & 0/8 \\
qr & 8/8 & 5/8 & 5/8 \\
spurious & 8/8 & 0/8 & 0/8 \\
collision & 0/8 & 8/8 & 8/8 \\
genuine & 8/8 & 8/8 & 8/8 \\
sibling & 0/4 & 1/4 & 1/4 \\
\bottomrule
\end{tabular}
\end{table}

\begin{table}[t]\centering\footnotesize
\caption{Registered predictions and outcomes (HITL: oracle).}
\label{tab:predictions}
\begin{tabular}{p{0.40\columnwidth}p{0.14\columnwidth}p{0.34\columnwidth}}
\toprule
Prediction & Outcome & Counts \\
\midrule
Clean packs: all arms tie & did not hold & baseline 8/8; agent\_rules 7/8; agent\_hybrid 7/8 \\
Collisions: baseline false-alarms on every one; agent at most 10\% & held & baseline 8/8, agent\_rules 0/8 \\
Misreads: baseline misses every flag; agent with oracle recovers at least 80\% at no more than 1 question per pack & did not hold & baseline detected 0/8; agent\_rules correct 6/8 at 0.75 questions per pack \\
Hybrid matches rules on verdicts & held & 60/60 identical labels \\
Hybrid beats rules on partial product matches & not evaluable & no item produced a partial product match \\
\bottomrule
\end{tabular}
\end{table}

}{\pendingbox{Tables pending: \texttt{paper/generated/results\_tables.tex} has not been written yet.}}
\else
\pendingbox{The evaluation tables (headline metrics with the oracle user and with questions skipped, per-scenario counts) are generated by \texttt{python cli.py eval} into \texttt{paper/generated/results.tex} and \texttt{results\_tables.tex}. They appear here once those files exist. No result in this report is typed by hand.}
\fi
